\documentstyle[%


righttag%


,11pt%


,amssymb%


,russian%               remove in Eng.


,izvru%                 Set izven% in Eng.


]{amsart}





%





\newcommand{\tts}[1]{}





\theoremstyle{plain}


\theorembodyfont{\it}


\newtheorem{thmnonum}{\hskip\parindent Теорема}


\renewcommand{\thethmnonum}{}


\renewcommand{\baselinestretch}{1.05}





%\hoffset=-15mm%                  For Eng.


\setcounter{page}{25}


\god{2004}


\nomer{1~(500)} %                        Remove in Eng.


%% remove ^^^^^^


%\nomer{Vol.\,48, No.\,1}%               Set in Eng.


%\str{\pageref{begin}--\pageref{end}}%  Set in Eng.


\udk{517.977}





\author{\it С.Е.\,Бухтояров, В.А.\,Емеличев}


% NB:   ^^ remove in Eng.


\title{О квазиустойчивости векторной траекторной задачи\\ с


параметрическим принципом оптимальности}


\thanks{Работа выполнена при финансовой поддержке


Белорусского государственного университета в рамках


Государственной программы фундаментальных исследований Республики


Беларусь ``Математические структуры 29''.}





\date{}


%\pagestyle{myheadings}%         Set in Eng.


\pagestyle{plain} %              Remove in Eng.


\begin{document}


%\thispagestyle{plain}%          Set in Eng.


\maketitle


\label{begin}








Под квазиустойчивостью векторной задачи дискретной оптимизации,


как обычно [1]--[5], будем понимать дискретный аналог


полунепрерывности снизу (по Хаусдорфу) многозначного отображения,


которое набору параметров задачи ставит в соответствие искомое


множество альтернатив (эффективных в том или ином смысле). Тем


самым квазиустойчивость задачи есть свойство сохранения


указанного множества при ``малых'' возмущениях параметров задачи,


а радиусом квазиустойчивости называется предельный уровень таких


возмущений.





В данной статье


рассматривается $n$-критериальная линейная комбинаторная


(на системе подмножеств конечного множества) задача


оптимизации, принцип оптимальности которой задается с помощью целочисленного


параметра $s$, изменяющегося в пределах от 1 до $n$. При этом


крайним значениям параметра соответствуют паретовский и слейтеровский


принципы оптимальности.


Для каждого значения параметра $s$ найдена формула радиуса


квазиустойчивости задачи, а также указаны необходимые и достаточные условия


устойчивости этого типа.








Класс векторных ($n$-критериальных) задач,


рассматриваемых в данной работе, можно описать следующей


комбинаторной моделью.


Пусть задан векторный критерий


$$


f(t,A)=(f_1(t,A_1),f_2(t,A_2),\dotsc,f_n(t,A_n)) \rightarrow


\min_{t \in T}


$$


с частными критериями вида


$$


f_i(t,A_i)=\sum_{j


\in N(t)}a_{ij},\quad i\in N_n,


$$


где $T$ --- множество траекторий,


$T\subseteq 2^E$, $|T| \geq 2$, $E=\{e_1,e_2,\dotsc,e_m\}$,


$m\geq 2$, $N_n=\{1,2,\dotsc,n\}$, $n\geq 1$, $N(t)=\{j\in N_n :


e_j\in t\}$, $A_i$ --- $i$-я строка матрицы $A=[a_{ij}]\in


\bold{R}^{n \times m}$. Будем полагать, что


$f_i(\emptyset,A_i)=0$.





С помощью бинарного отношения, зависящего от


целочисленного параметра, введем $n$ принципов


оптимальности (в иной терминологии --- функций выбора).


Для этого на множестве траекторий $T$ при всяком индексе


$s\in N_n$ введем бинарное отношение строгого предпочтения


$\Omega^n_s(A)$ между траекториями согласно формуле


$$


t \Omega^n_s(A) t^{\prime} \Leftrightarrow [t,t^{\prime},A]^+>


(s-1)[t,t^{\prime},A]^0+(n-1)[t,t^{\prime},A]^-,


$$


где


\begin{gather*}


[t,t^{\prime},A]^+ = |\{i \in N_n  : g_i(t,t^{\prime},A_i) >0\}|, \\


[t,t^{\prime},A]^0 = |\{i \in N_n :


g_i(t,t^{\prime},A_i) = 0\}|,\\


[t,t^{\prime},A]^- = |\{i \in N_n


: g_i(t,t^{\prime},A_i) < 0\}|,\\


g_i(t,t^{\prime},A_i)=f_i(t,A_i)-f_i(t^{\prime},A_i).


\end{gather*}





Это бинарное отношение порождает множество


$s$-эффективных траекторий


$$


T^n_s(A)=\{t \in T : T^n_s(t,A)=\emptyset\},


$$


где


$$


T^n_s(t,A)=\{t^{\prime}\in T : t \Omega^n_s(A) t^{\prime}\}.


$$





Легко видеть, что множество $T^n_1(A)$


совпадает с множеством Парето


$$


P^n(A)=\{t\in T : U^n(t,A)=\emptyset\},


$$


где


$$


U^n(t,A)=T^n_1(t,A)=\{t^{\prime}\in T :


[t,t^{\prime},A]^+>0\  \&\  [t,t^{\prime},A]^-=0\}.


$$


Очевидно, множество $T^n_n(A)$ есть


множество Слейтера (множество слабо эффективных


траекторий)


$$


S^n(A)=\{t\in T  :  V^n(t,A)=\emptyset\},


$$


где


$$


V^n(T,A)=T^n_n(t,A)=\{t^{\prime}\in T :[t,t^{\prime},A]^+=n\}.


$$





Отметим, что введенная здесь параметризация принципа


оптимальности ``от Парето до Слейтера'' является частным


случаем предложенного в [6] двупараметрического семейства $n^2$


принципов оптимальности в критериальном пространстве


$\bold{R}^n$.





Векторную задачу поиска множества $T^n_s(A)$ \ $s$-эффективных траекторий


будем обозначать через


$Z^n_s(A)$. В дальнейшем запись $t \Omega^n_s(A) t^{\prime}$


будем интерпретировать как утверждение о том, что траектория


$t^{\prime}$ доминирует траекторию $t$ в задаче $Z^n_s(A)$. Тем самым


траектория $t$ является $s$-эффективной тогда и только тогда, когда в


множестве $T$ не существует траекторий, которые доминируют ее в $Z^n_s(A)$.





Ясно, что $T^1_1(A)$ --- множество оптимальных траекторий обычной (скалярной)


линейной траекторной задачи, в схему которой вкладываются многие экстремальные


комбинаторные задачи на графах, а также задачи булева программирования и


ряд задач теории расписаний (см., напр., [7], [8]).





Следующие очевидные свойства вытекают из


приведенных выше определений и обозначений.





A1. Для любой матрицы $A\in\bold{R}^{n\times m}$ верны соотношения


$$


\emptyset\neq P^n(A)=T^n_1(A)\subseteq T^n_2(A) \subseteq \dotsb


\subseteq T^n_n(A)=S^n(A).


$$





A2.  Если $1 \leq s < k \leq n$, то для любой матрицы


$A\in \bold{R}^{n\times m}$ верна импликация


$t \Omega^n_k(A) t^{\prime} \Rightarrow t \Omega^n_s(A) t^{\prime}$.





A3. $[t,t^{\prime},A]^+=n\ \Leftrightarrow\ t \Omega^n_n(A) t^{\prime}$.





A4. $[t, t^{\prime},A]^+>0 \Rightarrow


t^{\prime} \overline{\Omega^n_1}(A) t$.





A5. Если $1\leq k < s\leq n$, то


$t^{\prime} \overline{\Omega^n_k}(A) t \Rightarrow t^{\prime}


\overline{\Omega^n_s}(A) t$.





\noindent{}Здесь и далее запись $t \overline{\Omega^n_s}(A) t^{\prime}$ будет


означать, что отношение


$t \Omega^n_s(A) t^{\prime}$ не выполняется, т.\,е. траектория $t^{\prime}$


не доминирует траекторию $t$ в задаче $Z^n_s(A)$.


Приведенные ниже очевидные свойства справедливы для


любого индекса $s\in N_n$.





A6. $(\forall t\in T\ (t^{\prime} \overline{\Omega^n_s}(A)\ t))


\Leftrightarrow t^{\prime}\in T^n_s(A)$.





A7. $(t^{\prime} \in T^n_s(A)\  \&\ t\in T\setminus \{t^{\prime}\})


\ \Rightarrow\ [t,t^{\prime},A]^++[t,t^{\prime},A]^0>0$.





Кроме того, нетрудно доказать [6], что бинарное отношение


$\Omega^n_s(A)$ при любом индексе\linebreak $s\in N_n$ антирефлексивно, асимметрично,


транзитивно и поэтому ациклично. Отметим, что отсюда следует непустота любого


множества $T^n_s(A)$, $s\in N_n$.





Как обычно [1]--[5], возмущение параметров векторного критерия


$f(t,A)$ будем осуществлять путем сложения матрицы $A\in


\bold{R}^{n\times m}$ с матрицами множества


$$


{\cal B}(\varepsilon)=\{B\in \bold{R}^{n\times m} :


\|B\|< \varepsilon\},


$$


где $\varepsilon >0$ --- предельный уровень


возмущений, $\|\cdot\|$ --- норма $l_{\infty}$ в пространстве


$\bold{R}^{n\times m}$, т.\,е.


$$


\|B\|=\max \{|b_{ij}| : (i,j)\in N_n\times N_m\},\quad


B=[b_{ij}]\in \bold{R}^{n\times m}.


$$





С учетом вышесказанного задача $Z^n_s(A)$ называется квазиустойчивой,


если


существует такое положительное число $\varepsilon>0$, что для любой


матрицы $B\in {\cal B}(\varepsilon)$ справедливо включение


$$


T^n_s(A) \subseteq T^n_s(A+B).


$$


Таким образом, задача $Z^n_s(A)$ квазиустойчива лишь в том случае, когда


при ``малых'' независимых возмущениях элементов


матрицы $A$ не исчезают


(хотя могут появляться новые)\linebreak


$s$-эффективные траектории исходной задачи.





Следуя [1], [3], [5], радиусом квазиустойчивости задачи $Z^n_s(A)$,


$s\in N_n$, назовем число


$$


\rho^n_s(A)=


\begin{cases}


\sup K(A), & \text{если}\ K(A) \neq \emptyset;\\ 0, & \text{если}\


K(A)=\emptyset,


\end{cases}


$$


где


$$


K(A)=\{\varepsilon >0  : \forall B\in{\cal B}(\varepsilon)


\ (T^n_s(A) \subseteq T^n_s(A+B))\}.


$$





Для траекторий $t, t^{\prime} \in T$ положим


$$


\Delta(t,t^{\prime})=|(t \cup t^{\prime})\setminus


(t \cap t^{\prime})|.


$$


Очевидно, $\Delta(t,t^{\prime}) \geq 1$ при $t \neq t^{\prime}$.





\begin{lem}


Если две различные траектории


$t,t^{\prime}\in T$ и матрица \mbox{$B\in \bold{R}^{n\times m}$}


таковы, что для некоторого индекса $k\in N_n$ выполняется


неравенство


\begin{equation} 


g_k(t,t^{\prime},A_k)>\|B_k\|\Delta(t,t^{\prime}),


\end{equation}


то для любого индекса $s\in N_n$ справедливо соотношение


$$


t^{\prime}\overline{\Omega^n_s}(A+B) t.


$$


\end{lem}





\begin{pf}


Поскольку очевидно неравенство


$$


g_k(t,t^{\prime},B_k) \geq -\|B_k\|\Delta(t,t^{\prime}),


$$


то с учетом (1) и линейности частных критериев выводим


$$


g_k(t,t^{\prime},A_k+B_k)=g_k(t,t^{\prime},A_k)+g_k(t,t^{\prime},B_k)


\geq g_k(t,t^{\prime},A_k)-\|B_k\|\Delta(t,t^{\prime})>0.


$$


Это значит, что $[t,t^{\prime},A+B]^+>0$. Поэтому согласно свойствам


A4 и A5 справедливо утверждение леммы 1.


\end{pf}





Пусть $0<\alpha < \varepsilon$, $t,t^{\prime}\in T$,


$t\neq t^{\prime}$. В дальнейшем будем использовать возмущающую


$n\times m$-мат\-ри\-цу $B^*(\alpha)=[b_{ij}^*]\in {\cal


B}(\varepsilon)$, элементы которой для любого индекса $i\in N_n$


задаются по формуле 


\begin{equation} 


b_{ij}^*=


\begin{cases}


-\alpha, & \text{если}\  j\in N(t\setminus t^{\prime});\\


 \alpha, & \text{если}\  j\in N(t^{\prime}\setminus t);\\


0 & \text{в остальных случаях.}


\end{cases}


\end{equation}


Очевидно, $\|B^*(\alpha)\|=\alpha$.





\begin{lem}


Пусть $t,t^{\prime}\in T$, $t\neq


t^{\prime}$, $s\in N_n$, $\alpha>0$. Если выполняются неравенства


\begin{equation} 


g_i(t,t^{\prime},A_i)<\alpha


\Delta(t,t^{\prime}),\ \ i\in N_n,


\end{equation}


то $$


t^{\prime} \Omega^n_s(A+B^*(\alpha)) t.


$$


\end{lem}








{\bf Доказательство.}


С учетом (3) и структуры возмущающей матрицы


$B^*(\alpha)\in {\cal B} (\varepsilon)$,


элементы которой задаются по формуле (2), легко убеждаемся в


справедливости следующих соотношений для любого индекса $i\in N_n$:


$$


g_i(t,t^{\prime},A_i+B^*_i(\alpha))=


g_i(t,t^{\prime},A_i)+g_i(t,t^{\prime},B^*_i(\alpha))=


g_i(t,t^{\prime},A_i)-\alpha\Delta(t,t^{\prime})<0.


$$


Отсюда получаем равенство $[t^{\prime},t,A+B^*(\alpha)]^+=n$


и потому в силу свойств A2 и A3 выводим


$$


t^{\prime} \Omega^n_s(A+B^*(\alpha)) t.\quad \square


$$





\begin{thmnonum}


При любых $n \geq 1$ и $s\in N_n$ радиус квазиустойчивости $\rho^n_s(A)$


векторной траекторной задачи $Z^n_s(A)$ имеет вид


\begin{equation} 


\rho^n_s(A)=\min_{t^{\prime}\in T^n_s(A)} \min_{t\in T\setminus


\{t^{\prime}\}}


\max_{i\in N_n} \frac{g_i(t,t^{\prime},A_i)}{\Delta(t,t^{\prime})}.


\end{equation}


\end{thmnonum}





\begin{pf}


Обозначим через $\varphi$ правую часть формулы (4).


Учитывая непустоту множества $T^n_s(A)$ и то, что для любой траектории


$t^{\prime}\in T^n_s(A)$ множество $T\setminus\{t^{\prime}\}$


также непусто, из


свойства A7 (ввиду $\Delta(t,t^{\prime})>0)$ выводим,


что число $\varphi$ существует и неотрицательно.





Далее докажем неравенство


\begin{equation} 


\rho^n_s(A) \geq \varphi.


\end{equation}


Будем полагать, что число $\varphi$ 


положительно (в случае, когда $\varphi=0$,


неравенство (5) очевидно). Пусть матрица


$B\in {\cal B}(\varphi)$. Тогда согласно


определению числа $\varphi$ для любых траекторий $t^{\prime}\in


T^n_s(A)$ и $t\in T\setminus\{t^{\prime}\}$ существует такой


индекс $k\in N_n$, при котором выполняется неравенство


$$


\|B_k\|\leq \|B\|<\varphi \leq \frac{g_k(t,t^{\prime},A_k)}


{\Delta(t,t^{\prime})}.


$$


Поэтому, применяя лемму 1, убеждаемся в


справедливости для любого индекса $s\in N_n$ соотношения


$$


t^{\prime} \overline{\Omega^n_s}(A+B) t.


$$


Отсюда, учитывая


антирефлексивность отношения $\Omega^n_s(\cdot)$, получаем


$$


\forall


B\in {\cal B}(\varphi)\ \ \forall t^{\prime}\in T^n_s(A)\ \ \forall


t\in T\ \ (t^{\prime} \overline{\Omega^n_s}(A+B) t).


$$


Поэтому


согласно свойству A6 траектория $t^{\prime}$ является


$s$-эффективной траекторией в любой возмущенной задаче


$Z^n_s(A+B)$, $B\in {\cal B}(\varphi)$. Это значит, что для любой


матрицы $ B\in {\cal B}(\varphi)$ справедливо включение


$$


T^n_s(A+B)\supseteq T^n_s(A).


$$


Следовательно, верно неравенство (5).





Наконец, докажем неравенство


\begin{equation} 


\rho^n_s(A) \leq \varphi.


\end{equation}





Пусть $\varphi < \alpha < \varepsilon$. Тогда согласно определению


числа $\varphi$ существуют такие траектории\linebreak $t^{\prime}\in


T^n_s(A)$ и $t\in T\setminus \{t^{\prime}\}$, что для любого


индекса $i\in N_n$ верны неравенства


$$


g_i(t,t^{\prime},A_i)\leq


\varphi\Delta(t,t^{\prime})< \alpha\Delta(t,t^{\prime}).


$$


Отсюда, пользуясь леммой 2, получаем


$$


t^{\prime} \Omega^n_s(A+B^*(\alpha)) t.


$$


Из этого соотношения


следует, что $t^{\prime}$ не является $s$-эффективной траекторией


в возмущенной задаче $Z^n_s(A+B^*(\alpha))$. В результате выводим


$$


\forall \varepsilon > \varphi\ \ \exists B^*(\alpha)\in {\cal


B}(\varepsilon) \ \ (T^n_s(A+B^*(\alpha))\not \supseteq T^n_s(A)).


$$


Это свидетельствует о справедливости неравенства (6), которое


вместе с неравенством (5) дает формулу (4).


\end{pf}





В принятых обозначениях традиционное множество Смейла, т.\,е. множество


строго эффективных траекторий, задается следующим образом:


$$


Q^n(A)=\{t\in T : W(t,A)=\emptyset\},


$$


где


$$


W(t,A)=\{t^{\prime}\in T\setminus \{t\} : [t,t^{\prime},A]^-=0\}.


$$


Очевидно, что $Q^n(A)\subseteq P^n(A)$, при этом множество $Q^n(A)$ может


быть пустым.





\begin{cor}


Для любых индексов $s \in N_n$ и $n\geq 1$


эквивалентны утверждения





--- задача $Z^n_s(A)$ квазиустойчива,





--- $T^n_s(A)=Q^n(A)$,





--- $\forall t^{\prime}\in T^n_s(A)$, $\forall t\in


T\setminus \{t^{\prime}\}$ $([t,t^{\prime},A]^+>0)$.


\end{cor}





Отсюда, в частности, вытекает известный результат [7]: скалярная задача


$Z^1_1(A)$ квазиустойчива тогда и только тогда, когда


имеет единственную оптимальную траекторию.





Ввиду свойства A1 из теоремы также вытекает





\begin{cor}


Для любой матрицы $A\in \bold{R}^{n\times m}$,


$n, m \geq 2$, верны неравенства


$$


\rho^n_1(A)\geq \rho^n_2(A)\geq \dotsb\geq \rho^n_n(A).


$$


\end{cor}





Отсюда легко выводим





\begin{cor}


Если $1\leq s < k\leq n$, $n\geq 2$, то


из квазиустойчивости задачи $Z^n_k(A)$ следует квазиустойчивость


задачи $Z^n_s(A)$.


\end{cor}





В заключение заметим, что при $s=1$ теорема


и следствие 1 переходят в результаты, полученные в [1] и


касающиеся квазиустойчивости векторной траекторной задачи


$Z^n_1(A)$ поиска множества Парето.








\begin{thebibliography}{9}








\bibitem{1}


Емеличев В.А., Кравцов М.К., Подкопаев Д.П. {\em О


квазиустойчивости траекторных задач векторной оптимизации} //


Матем. заметки. -- 1998. -- Т.\,63.


-- Вып.\,1. -- С.\,21--27.





\bibitem{2}


Емеличев В.А., Степанишина Ю.В. {\em Квазиустойчивость


векторной нелинейной траекторной задачи с паретовским принципом


оптимальности} // Изв. вузов. Математика. --


2000. -- \No\,12. -- С.\,27--32.





\bibitem{3}


Емеличев В.А., Подкопаев Д.П. {\em Устойчивость и


регуляризация векторных задач целочисленного линейного


программирования} // Дискрет. анализ и исследов. операций. Сер.\,2.


-- 2001. -- Т.\,8. -- \No\,1. -- С.\,47--69.





\bibitem{4} Емеличев В.А., Степанишина Ю.В. {\em О квазиустойчивости


векторной траекторной задачи мажоритарной оптимизации} // Матем.


заметки. -- 2002. -- Т.\,72. -- \No\,1. -- С.\,38--47.





\bibitem{5} Emelichev V.A., Girlich E., Nikulin Yu.V., Podkopaev D.P.


{\em Stability and regularization of vector problems of integer


linear programming} // Optimization. -- 2002. -- V.\,51. -- \No\,4. --


P.\,645--676.





\bibitem{6} Емеличев В.А., Пашкевич А.В. {\em О параметризации


принципа оптимальности в критериальном пространстве} // Дискрет.


анализ и исследов. операций. Сер.\,2. -- 2002. -- Т.\,9. -- \No\,1.


-- С.\,21--32.





\bibitem{7} Леонтьев В.К. {\em Устойчивость в линейных дискретных


задачах}


// Пробл. кибернетики. -- М.: Наука, 1979. -- Вып.\,35. -- С.\,169--184.





\bibitem{8} Sotskov Yu.N., Leontev V.K., Gordeev E.N. {\em Some


concepts of stability analysis in combinatorial optimization} //


Discrete Appl. Math. -- 1995. -- V.\,58. -- \No\,2. -- P.\,169--190.





\end{thebibliography}





\vskip 0.5cm





\post{\it Белорусский государственный}{\it Поступила}


\post{\it университет}{03.02.2003}





\label{end}


\end{document}


